% SPDX-License-Identifier: LPPL-1.3c
% Copyright (C) 2026 Christophe Jorssen
% Author and Current Maintainer: Christophe Jorssen <christophe.jorssen@gmail.com>
% LPPL maintenance status: maintained
% This file is part of luafemm. See LICENSE and LICENSES.md for its terms.
\section{Complete examples in the three formats}
\label{sec:format-examples}
The same magnetic scene can be compiled with plain LuaTeX, LuaLaTeX and
ConTeXt MkIV. The following listings show the shared geometry and every
complete wrapper, followed by its actual output. Place \texttt{scene.tex}
beside the selected wrapper and compile from that directory with an
installed luafemm TEXMF tree. The development build does this automatically.
The geometry, profile and solver are identical; page layout follows the format.
The manual crops the empty paper margins around each output for readability.

\subsection{A minimal plain-LuaTeX document}
This self-contained example uses the same picture keys and physical paths as
LaTeX and ConTeXt. It needs no format-specific wrapper or older command interface.
\luafemmexample{plain}

\subsection{Shared physical scene}
The alnico rectangle, iron return piece and passive copper region precede
the first field request. The curved measurement path is registered before
the field is solved. Each wrapper subsequently plots that same profile.
\luafemmsource{formats/scene}

\subsection{Plain LuaTeX}
Compile with \texttt{luatex plain.tex}. Plain environment commands are
used directly; the source requires no LaTeX wrapper.
\luafemmexample{formats/plain}

\subsection{LuaLaTeX}
Compile with \texttt{lualatex latex.tex}. The package wrapper and the
library entry point lead to the same generic implementation.
\luafemmexample{formats/latex}

\subsection{ConTeXt MkIV}
Compile with \texttt{context --luatex context.tex}. The explicit engine
selection matters on installations whose default ConTeXt engine is LMTX.
\luafemmexample{formats/context}
